\cotefile{157-2}
\scanpage{338}
\mspage{77}
% archive sheet number: 337
À la page 61, j'ai invoqué l'existence d'un \uline{foncteur} de ``\uline{composition des chemins}''
\[
\underline{\mathrm{Ch}}_{\Sigma}(X;a,b)\times\underline{\mathrm{Ch}}_{\Sigma}(X;b,c)\mathrel{\unsure{-|\longrightarrow}}\underline{\mathrm{Ch}}_{\Sigma}(X;a,c),
\]
qui sur les fibres du premier \unsure{mb} au dessus de $(\tau,\tau')\in\mathrm{Ob}(\underline{\mathrm{Ty}}\times\underline{\mathrm{Ty}})$, serait donné par
\[
\underline{\mathrm{Ch}}_{\Sigma,\tau}(X;a,b)\times\underline{\mathrm{Ch}}_{\Sigma,\tau'}(X;b,c)\longrightarrow\underline{\mathrm{Ch}}_{\Sigma,\tau\tau'}(X;a,c),\tag*{$(\unsure{*})$}
\]
la composition des chemins de types précisés $\tau,\tau'$. Qu'on ait là un \uline{foncteur}, pas seulement une application sur les ensembles d'objets, est assez évident, et peut s'exprimer par exemple à l'aide du diagramme d'amalgamation
\diagram{\begin{tikzcd}[ampersand replacement=\&,row sep=small,column sep=small]
\& \unsure{\varepsilon|E}=\{0,1\}\arrow[dl,hook]\arrow[dr,hook'] \&\\
I_{\tau}\arrow[dr,hook] \& \& I_{\tau'}\arrow[dl,hook',no head]\\
\& I_{\tau\tau'} \&
\end{tikzcd}}
en le transformant par $\underline{\mathrm{Hom}}(\bullet,X)$ (au lieu de lui \deleted{\unsure{appliquer}} transformer par $\mathrm{Hom}_{\mathrm{Cat}}(\cdot,X)$ seulement), \deleted{\unsure{bien sûr pour la catégorie}} \inserted{du moins dans le cas} particulier où $\Sigma=\mathrm{Fl}(X)$ (qui était celui de la p. 61). Mais le passage aux flèches est assez évident -- pour deux composés
\diagram{\begin{tikzcd}[ampersand replacement=\&,column sep=large,row sep=small]
a\arrow[r,"c_1",""{name=A,below}]\arrow[d,phantom,"\shortparallel" description] \& b\arrow[r,"{c'_1}",""{name=B,below},dash]\arrow[d,phantom,"\shortparallel" description] \& c\arrow[d,phantom,"\shortparallel" description]\\
a\arrow[r,"c_2",""{name=C,above}] \& b\arrow[r,"{c'_2}",""{name=D,above}] \& c
\arrow[from=A,to=C,Rightarrow,"{\unsure{w|u}}"']
\arrow[from=B,to=D,Rightarrow,"{w'}"]
\end{tikzcd}}
$c'_1\circ c_1$ et $c'_2\circ c_2$, de sorte que les hom \inserted{$\unsure{w''|w^{4}}$} du premier dans le \unsure{second}, dans $\underline{\mathrm{Ch}}_{\Sigma}(X;a,c)$, \deleted{``\unsure{sont exactement les couples}} \inserted{\unsure{qui|font} \unsure{sont|sur} l'\unsure{identité|identific$^{\mathrm{tion}}$} sur le sommet} d'indice $n+1$ des \inserted{deux} chemins (où $n=\mathrm{long}\,\tau$), \unsure{correspondant|correspond} à \deleted{\unsure{$\mathrm{id}_{\Sigma,\tau}$}} \unsure{$1_{\tau}$} et à \unsure{$0_{\tau'}$}, ``sont'' exactement les couples
